\documentclass[11pt]{article}
\usepackage[margin=1in]{geometry}
\usepackage{amsmath,amssymb,amsthm}
\usepackage{booktabs}
\usepackage{hyperref}
\usepackage{enumitem}
\hypersetup{colorlinks=true,linkcolor=black,urlcolor=blue,citecolor=black}
\title{Hollow Purple Semantic Nullification\\\large Contradictory Semantic Composition, Uncertainty, and Representational Stress in Large Language Models}
\author{Aryan Singh Chandel\\\small ORCID: \href{https://orcid.org/0009-0007-1774-2480}{0009-0007-1774-2480}}
\date{\today}

\begin{document}
\maketitle

\begin{abstract}
We study whether tightly coupled semantic contradictions induce a measurable instability regime in large language models. We call this regime semantic nullification: not literal model failure, but a reproducible reduction in representational coherence under mutually exclusive semantic pressure. Building on prior work in antonymy, contradiction benchmarks, uncertainty estimation, and interpretability, we define a mathematical framework for quantifying contradiction-induced stress through token entropy, semantic entropy, hidden-state variance, and contradiction-resolution failure. The goal is to convert the original conjecture about antonym collision in shared semantic space into a precise empirical hypothesis.
\end{abstract}

\section{Introduction}
Contradiction is often evaluated behaviorally, especially in natural language inference. This paper asks a different question: what happens when contradictory meanings are forced into a shared prompt state during generation? We hypothesize that contradiction can act as a controlled semantic stressor and produce elevated uncertainty, reduced hidden-state stability, and lower output coherence.

\section{Mathematical Formulation}
Let $x$ denote a prompt, $M$ a language model, and $h_\ell(x)$ the hidden state of $x$ at layer $\ell$.

\subsection{Token Entropy}
For decoding step $t$, define
\[
H_t(x) = - \sum_i p_M(i \mid x, y_{<t}) \log p_M(i \mid x, y_{<t}).
\]
The average decoding entropy over $T$ steps is
\[
H_{\mathrm{avg}}(x) = \frac{1}{T} \sum_{t=1}^T H_t(x).
\]

\subsection{Hidden-State Drift}
Given a control prompt $x_c$ and contradiction prompt $x_k$, define
\[
D_\ell(x_c, x_k) = 1 - \cos\big(h_\ell(x_c), h_\ell(x_k)\big).
\]

\subsection{Semantic Variance}
Given repeated sampled outputs clustered into semantic groups with masses $q_j$, define
\[
SE(x) = -\sum_j q_j \log q_j.
\]

\subsection{Nullification Index}
We define a composite study index
\[
NI(x) = \alpha H_{\mathrm{avg}}(x) + \beta SE(x) + \gamma V_\ell(P_x) + \delta R(x),
\]
where $V_\ell(P_x)$ denotes paraphrase-level hidden-state variance and $R(x)$ denotes contradiction-resolution failure.

\section{Experimental Design}
We compare coherent controls, mild tension prompts, explicit contradictions, and deep contradictions. The empirical benchmark in \texttt{benchmark\_suite.json} and \texttt{run\_benchmark.py} provides the first-pass measurement scaffold.

\section{Assumptions}
This study assumes that contradiction depth is experimentally manipulable, that hidden-state and uncertainty metrics are sensitive to contradiction-induced stress, and that matched ambiguity controls can separate contradiction effects from generic prompt difficulty.

\section{Formal Claim Structure}

\subsection{Definition}
Semantic nullification is the condition in which contradiction-bearing prompts exhibit measurable instability relative to matched controls in uncertainty, semantic divergence, or representational stability.

\subsection{Proposition H1}
Increasing contradiction depth increases average uncertainty.

\subsection{Test Statistic for H1}
Compare condition-wise means of $H_{\mathrm{avg}}(x)$ across contradiction levels and ambiguity controls.

\subsection{Decision Rule for H1}
Support H1 only if contradiction conditions show a stable directional increase in $H_{\mathrm{avg}}(x)$ beyond matched ambiguity controls across repeated runs.

\subsection{Failure Mode for H1}
H1 fails if uncertainty rises equally under non-contradictory ambiguity, or if no monotonic trend appears with contradiction depth.

\subsection{Proposition H2}
Increasing contradiction depth increases semantic output divergence.

\subsection{Test Statistic for H2}
Compute semantic entropy $SE(x)$ or repeated-output cluster dispersion for each condition.

\subsection{Decision Rule for H2}
Support H2 only if contradiction conditions produce consistently higher semantic divergence than controls under fixed decoding settings.

\subsection{Failure Mode for H2}
H2 fails if repeated outputs remain as stable as control outputs or if divergence is explained entirely by decoding randomness.

\subsection{Proposition H3}
Increasing contradiction depth decreases representational stability.

\subsection{Test Statistic for H3}
Measure $D_\ell(x_c,x_k)$ and within-family hidden-state variance $V_\ell(P_x)$ across layers.

\subsection{Decision Rule for H3}
Support H3 only if contradiction families show persistently larger drift or variance than coherent controls across relevant layers.

\subsection{Failure Mode for H3}
H3 fails if internal geometry remains indistinguishable from controls or if instability is confined to output text alone.

\section{Validation}
The theory is supported if contradiction depth predicts increases in uncertainty and instability beyond ambiguity and prompt difficulty controls. We require convergent signals across entropy, variance, and evaluator-rated coherence.

\section{Falsifiers}
The theory is weakened if contradiction prompts fail to separate from ambiguity controls, if no monotonic relationship appears across contradiction depth, or if the effect vanishes across models and paraphrases.

\section{Evidence and Proof Standard}
In this project, evidence means stable metric separation under controls. Strong evidence requires replication across models, paraphrases, and multiple measurement channels. Proof, in the empirical sense, requires that the nullification hypothesis survive these tests without a simpler competing explanation within scope.

\section{Conclusion}
Semantic nullification is not yet proven, but it is mathematically specifiable and empirically testable.

\nocite{*}
\bibliographystyle{plain}
\bibliography{references}
\end{document}

